\textbf{Время:} резерв

\textbf{Тезис:} Rlong добавляет к поперечным радиусам характеристику области испускания вдоль пучка.

\textbf{На экране:} название Rlong и оси увеличенного графика; правая легенда тех же центральностей; сначала красная последовательность по kT, затем сопоставление красных и синих точек при большом kT.

\textbf{Текст:} Это третий параметр того же анализа — радиус long, вдоль пучка. График увеличен, но горизонтальная ось осталась прежней: k-T. По вертикали снова радиус в фемтометрах.

Справа повторена та же цветовая легенда центральностей, что на предыдущем экране. Проследим красные точки: с ростом k-T продольный радиус в этом архивном результате заметно уменьшается. Синие точки также идут вниз, а при больших k-T красная и синяя последовательности сближаются.

Таким образом, зависимость long имеет свою форму. Одним общим размером заменить три направления нельзя: out, side и long показывают разные характеристики области испускания.

Теперь у нас есть три набора модельных радиусов, и у STAR измерены параметры с теми же названиями. Их можно поставить рядом, сохраняя различия условий анализа.

\textbf{Переход:} Перенесём радиусы на общие оси и посмотрим, какие выводы допускает такое сравнение.

\textbf{Справка:} Балинов (2025), исходный слайд 17, нижняя левая панель. Это параметр исходной трёхмерной аппроксимации, не независимое измерение времени жизни и не размер всего источника. Моменты Rlong зависят от выбранных пар и системы отсчёта. Исходные точки и погрешности сохранены. Реализация аппроксимации 2025 года и физическая корректность её численного масштаба независимо не проверены; точные количественные выводы по этому архиву не заявляются.
